\documentclass[11pt,oneside,openany]{book}
\usepackage[paperwidth=6.5in,paperheight=9.5in,inner=.72in,outer=.72in,top=.75in,bottom=.78in,headheight=14pt]{geometry}
\usepackage{fontspec}
\setmainfont{lmroman10-regular.otf}[BoldFont=lmroman10-bold.otf,ItalicFont=lmroman10-italic.otf,BoldItalicFont=lmroman10-bolditalic.otf,SmallCapsFont=lmromancaps10-regular.otf]
\usepackage[main=ngerman]{babel}
\usepackage{amsmath,amssymb,graphicx,microtype,fancyhdr,needspace,tikz}
\usepackage[unicode,hidelinks]{hyperref}
\hypersetup{pdftitle={Carl Friedrich Gauss: Werke, Band V. Anziehung, Mechanik, Flüssigkeitsgleichgewicht, Erdmagnetismus und Allgemeine Lehrsätze, Dioptrische Untersuchungen, Anzeigen eigener Abhandlungen, Einleitung zu den Beobachtungen des Magnetischen Vereins von 1836, Ein neues Hilfsmittel für die magnetischen Beobachtungen, Ein neues Instrument zur Messung von Intensitätsänderungen, Anleitung zur Bestimmung der Schwingungsdauer einer Magnetnadel, Über ein Mittel, die Beobachtung von Ablenkungen zu erleichtern, Zur Bestimmung der Konstanten des Bifilarmagnetometers, Vorschriften zur Bestimmung der magnetischen Wirkung, die ein Magnetstab in der Ferne ausübt, Über die Anwendung des Magnetometers zur Bestimmung der absoluten Deklination, Beobachtungen der magnetischen Inklination in Göttingen, Fundamentalgleichungen für die Bewegung schwerer Körper auf der rotierenden Erde, Über achromatische Doppelobjektive, Brief von Gauss an Brandes, Berichtigung der Schneiden einer Waage, Physikalische Beobachtungen, Nordlicht am 7. Januar 1831, Magnetisches Observatorium in Göttingen (August 1834), Beobachtungen der magnetischen Variation in Göttingen und Leipzig, Magnetische Beobachtungen (März 1835), Magnetnadel in Kopenhagen und Mailand, Brief von Gauss an Schumacher (1836), Verfahren bei den Beobachtungsterminen, Dreijährige tägliche Deklinationsbeobachtungen, Terminszeichnungen und Beobachtungszahlen von 1836, Terminszeichnungen und Beobachtungszahlen von 1837, Zu Wilkes’ Beobachtungen in der Nähe des magnetischen Südpols, Besprechung von Benzenbergs Über die Daltonsche Theorie, Besprechung von Fischers Anleitung zur Herstellung künstlicher Magnete, Besprechung der Ergebnisse des Magnetischen Vereins für 1836, Das mathematisch-physikalische Institut in Marburg, Zur mathematischen Theorie der elektrodynamischen Wirkungen, §§ 1–21, Kugelfunktionen, Zum Gebrauch des Komparators, Wirkung eines leuchtenden Punkts, Scherings Bemerkungen, Historisches Inhaltsverzeichnis und Druckvermerk},pdfauthor={Carl Friedrich Gauss}}
\setlength{\parindent}{1.2em}\setlength{\parskip}{0pt plus .3pt}
\setlength{\emergencystretch}{1.4em}
\pagestyle{fancy}\fancyhf{}
\fancyhead[L]{\small THEORIE DER ANZIEHUNG}
\fancyfoot[C]{\thepage}\renewcommand{\headrulewidth}{.25pt}
\newcommand{\SourceLeaf}[1]{\hypertarget{native-#1}{}\label{native-#1}}
\newcommand{\Article}[1]{\par\medskip\Needspace{4\baselineskip}\pdfbookmark[1]{Art. #1}{attractionis-art-#1}\centerline{\textbf{#1.}}\nobreak\smallskip}
\renewcommand{\thefootnote}{*)}
\newcommand{\HistoricalNote}[1]{\footnote{\csname HistoricalNoteText#1\endcsname}}
\newcommand{\BareRoot}[2][]{\mathord{{}^{#1}\!\surd}#2}
\widowpenalty=10000\clubpenalty=10000\displaywidowpenalty=10000
\raggedbottom
